\section{Desarrollo}
\[
Z_{9602} =
\begin{bmatrix}200.0\angle 0.0^\circ\,\Omega & 215.3\angle -4.0^\circ\,\Omega \\ 1415\angle -58.0^\circ\,\Omega & 431.1\angle -4.0^\circ\,\Omega\end{bmatrix}
\qquad
Z_{9610} =
\begin{bmatrix}289.7\angle -2.0^\circ\,\Omega & 215.2\angle -27.0^\circ\,\Omega \\ 159.7\angle -36.0^\circ\,\Omega & 301.5\angle -2.0^\circ\,\Omega\end{bmatrix}
\]

\subsection{Matrices $Y$ experimentales}
Los valores se expresan en mS.
\[
Y_{9602,\mathrm{exp}} =
\begin{bmatrix}0.470\angle 0.0^\circ\,\mathrm{{mS}} & 0.240\angle -7.7^\circ\,\mathrm{{mS}} \\ 2.441\angle -4.0^\circ\,\mathrm{{mS}} & 0.361\angle -4.0^\circ\,\mathrm{{mS}}\end{bmatrix}
\qquad
Y_{9610,\mathrm{exp}} =
\begin{bmatrix}3.704\angle 0.0^\circ\,\mathrm{{mS}} & 0.273\angle 0.0^\circ\,\mathrm{{mS}} \\ 2.593\angle 162.0^\circ\,\mathrm{{mS}} & 0.456\angle 0.0^\circ\,\mathrm{{mS}}\end{bmatrix}
\]

\subsection{Matrices $T$ experimentales}
\[
T_{9602,\mathrm{exp}} =
\begin{bmatrix}0.141\angle 58.0^\circ\, & 409.6\angle -176.0^\circ\,\Omega \\ 0.707\angle 58.0^\circ\,\mathrm{mS} & 0.193\angle -176.0^\circ\,\end{bmatrix}
\qquad
T_{9610,\mathrm{exp}} =
\begin{bmatrix}1.814\angle 34.0^\circ\, & 385.7\angle 18.0^\circ\,\Omega \\ 6.263\angle 36.0^\circ\,\mathrm{mS} & 1.429\angle 18.0^\circ\,\end{bmatrix}
\]

\subsection{Verificación por transformación desde $Z$}
Para comparar parámetros se usaron:
\[
Y = Z^{-1},
\qquad
T_Z =
\begin{bmatrix}
\dfrac{Z_{11}}{Z_{21}} &
\dfrac{Z_{11}Z_{22}-Z_{12}Z_{21}}{Z_{21}} \\[1.2em]
\dfrac{1}{Z_{21}} &
\dfrac{Z_{22}}{Z_{21}}
\end{bmatrix}.
\]

\begin{table}[H]
\centering
\caption{Comparación de $Y_\mathrm{exp}$ contra $Y=Z^{-1}$. Se informa error relativo por entrada.}
\begin{tabular}{c c}
\toprule
Cuadripolo & Error relativo de $Y$ \\
\midrule
9602 & $\begin{bmatrix}111.8\% & 100.1\% \\ 60.5\% & 117.3\%\end{bmatrix}$ \\
9610 & $\begin{bmatrix}35.7\% & 106.9\% \\ 73.8\% & 89.0\%\end{bmatrix}$ \\
\bottomrule
\end{tabular}
\end{table}

\begin{table}[H]
\centering
\caption{Comparación de $T_\mathrm{exp}$ contra $T$ obtenido desde $Z$. Se informa error relativo por entrada.}
\begin{tabular}{c c}
\toprule
Cuadripolo & Error relativo de $T$ \\
\midrule
9602 & $\begin{bmatrix}0.0\% & 130.3\% \\ 0.0\% & 148.7\%\end{bmatrix}$ \\
9610 & $\begin{bmatrix}0.0\% & 60.1\% \\ 0.0\% & 34.3\%\end{bmatrix}$ \\
\bottomrule
\end{tabular}
\end{table}

\textbf{Observación importante:} los parámetros $A$ y $C$ de $T$ coinciden prácticamente de forma exacta con los calculados a partir de $Z$, pero $B$ y $D$ no. 

\section{Conexión en cascada}
Para una conexión en cascada se espera:
\[
T_{eq,esp} = T_{9602}\,T_{9610}.
\]

\[
T_{eq,\mathrm{exp}} =
\begin{bmatrix}7.213\angle 81.0^\circ\, & 1368\angle 56.0^\circ\,\Omega \\ 31.48\angle 81.0^\circ\,\mathrm{mS} & 5.789\angle 56.0^\circ\,\end{bmatrix}
\]
\[
T_{eq,\mathrm{esp}} =
\begin{bmatrix}2.416\angle -144.8^\circ\, & 554.9\angle -162.6^\circ\,\Omega \\ 1.093\angle 152.4^\circ\,\mathrm{mS} & 0.249\angle 139.5^\circ\,\end{bmatrix}
\]

\begin{table}[H]
\centering
\caption{Error relativo entre cascada experimental y cascada esperada a partir de los cuadripolos individuales.}
\begin{tabular}{c}
\toprule
Error relativo de $T_{eq}$ \\
\midrule
$\begin{bmatrix}375.1\% & 330.7\% \\ 2850.5\% & 2319.5\%\end{bmatrix}$ \\
\bottomrule
\end{tabular}
\end{table}

ESTO DA COMO EL CULOOOOOOOOOO (NO ENTREGAR ESTO)

\section{Conexión en serie}
Según la planilla, se conectó el cuadripolo 9602 en serie con el 9610 invertido. Si se cumple la condición de Brune, se espera:
\[
Z_{eq,esp} = Z_{9602} + Z_{9610,\mathrm{invertido}}.
\]

\[
Z_{eq,\mathrm{exp}} =
\begin{bmatrix}687.7\angle 0.0^\circ\,\Omega & 492.6\angle -1.4^\circ\,\Omega \\ 92.05\angle 0.5^\circ\,\Omega & 819.1\angle 0.0^\circ\,\Omega\end{bmatrix}
\]
\[
Z_{eq,\mathrm{esp}} =
\begin{bmatrix}501.4\angle -1.2^\circ\,\Omega & 360.7\angle -17.6^\circ\,\Omega \\ 1603\angle -54.0^\circ\,\Omega & 720.6\angle -3.2^\circ\,\Omega\end{bmatrix}
\]

\begin{table}[H]
\centering
\caption{Error relativo entre conexión serie experimental y conexión serie esperada.}
\begin{tabular}{c}
\toprule
Error relativo de $Z_{eq}$ \\
\midrule
$\begin{bmatrix}37.2\% & 49.1\% \\ 96.8\% & 14.9\%\end{bmatrix}$ \\
\bottomrule
\end{tabular}
\end{table}

En particular, $Z_{21}$ presenta una diferencia muy grande. Esto sugiere que la condición de Brune no quedó garantizada.
\newpage